\chapter{Conclusiones}\label{cap9}

El desarrollo del presente Trabajo Terminal 1 finalizó con la obtención de los
parámetros para el diseño detallado del sistema compuesto por el biorreactor
(tanque de pruebas) y unidad de control (interfaz HMI), así como la validación
teórica de la viabilidad de su construcción.

En el ámbito de la mecánica, la geometría del tanque fue validada mediante análisis
por elementos finitos, simulaciones de presión hidrostática y manométrica fueron
efectuadas y concluidas con éxito asegurando un factor de seguridad superior a 7,
lo que permite una operación segura. Se consideraron algunos factores de diseño
presentados en la bibliografía adaptados a las prioridades de este prototipo, pero
sin comprometer su operación.

Se abordó también el diseño del sistema de control, donde se establecieron
diferentes técnicas dependiendo de la naturaleza de la variable, tanto la lógica
difusa como el control proporcional integral fueron validados teóricamente mediante
simulaciones para asegurar que el sistema alcanzase los estados deseados.

Finalmente se realizó un análisis de costos, así como el plan de manufactura con
el objetivo de comprobar la viabilidad tanto económica como técnica del prototipo,
arrojando un presupuesto competitivo frente a soluciones comerciales sin
personalización para este tipo específico de hongos.

En conclusión, el diseño planteado a lo largo de este documento satisface los
requerimientos biológicos planteados para el cultivo de los ascomicetos mediante
una solución mecatrónica viable y adecuada a la problemática. Además, abre el
camino para la futura etapa de construcción física, integración electrónica y pruebas
a realizar a lo largo de Trabajo Terminal II.
